\chapter{Manual d'operació i procediments}
\label{chap:manual}

Aquest capítol recull la configuració de referència del sistema i el procediment de treball, aplicables a qualsevol vol de fotogrametria. Els resultats obtinguts amb aquests procediments es presenten al \cref{chap:vols-prova}.

\section{Antenes GNSS}
\label{sec:antenes-gnss}

\subsection{Antenes helicals del dron}

Les dues antenes helicals del \emph{rover} van muntades:
\begin{itemize}
    \item Verticals, paral·leles entre elles, apuntant amunt.
    \item Al punt més alt del dron.
    \item Amb una separació mínima de 10~cm respecte al VTx, els ESC i els cables de potència, per evitar interferències.
    \item Sense inclinació respecte a l'eix vertical.
\end{itemize}

\subsection{\emph{Baseline} del \emph{moving baseline}}

La separació entre les dues antenes helicals és de \textbf{0,38~m}, sobre un \emph{bracket} rígid de carboni alineat amb l'eix Y (lateral) del dron. Els \emph{offsets} respecte al centre de gravetat són $Y = \mp 0{,}19$~m i $Z = -0{,}10$~m (GPS1 a l'esquerra, GPS2 a la dreta), amb $X = 0$ (\cref{tab:params-gps-mb}).

La precisió de \emph{heading} d'un \emph{moving baseline} escala de manera inversament proporcional a la longitud del \emph{baseline}: amb 0,38~m s'estima una $\sigma$ de l'ordre de $0{,}5^{\circ}$, que cal validar de manera empírica (\cref{sec:validacio}). Qualsevol canvi en la longitud del \emph{bracket} requereix recalibrar el \emph{heading} GNSS.

\subsection{Antena de la base}

Antena de tipus \emph{patch} horitzontal amb pla de massa (\emph{ground plane}), anivellada sobre el trípode topogràfic.

\subsection{Sincronisme de la càmera}

El pin central de la sabata (\emph{hot-shoe}) de la càmera és el senyal \textbf{X-sync}: el Seagull SYNC2 hi encaixa directament, ja que la Sony $\alpha$6000 disposa d'una \emph{Multi Interface Shoe} compatible amb l'estàndard ISO~518.

\section{Configuració de ports sèrie del FC (Kakute H7)}
\label{sec:ports-serie}

\begin{table}[H]
    \centering
    \caption{Assignació de ports sèrie del Kakute H7}
    \label{tab:ports-serie}
    \begin{tabular}{lllc}
        \toprule
        Port & Funció & Baud rate & Protocol \\
        \midrule
        SERIAL0 (USB) & MAVLink2 (configuració) & -- & -- \\
        SERIAL1 (UART1, port VTX) & RFD868x MAVLink & 57600 & 2 \\
        SERIAL2 (UART2) & Bluetooth intern (ocupat) & -- & Deshabilitat \\
        SERIAL3 (UART3) & F9P \#1 -- GPS1 (MB-Base) & 230400 & 5 \\
        SERIAL4 (UART4) & F9P \#2 -- GPS2 (MB-Rover) & 230400 & 5 \\
        SERIAL6 (UART6) & RC ELRS (CRSF) & 420000 & 23 \\
        SERIAL7 (UART7, port ESC) & Telemetria ESC & 115200 & 16 \\
        \bottomrule
    \end{tabular}
\end{table}

\section{Configuració de la càmera}
\label{sec:config-camera}

\subsection{Paràmetres del controlador}

A partir d'ArduPilot 4.5, els paràmetres de càmera tenen el prefix \texttt{CAM1\_*} per donar suport multicàmera; la documentació i els fòrums antics encara fan servir noms \texttt{CAM\_*}.

\begin{table}[H]
    \centering
    \caption{Paràmetres de càmera (ArduPilot 4.5 o superior)}
    \label{tab:params-camera}
    \begin{tabular}{lP{9.5cm}}
        \toprule
        Paràmetre & Valor / Funció \\
        \midrule
        \texttt{SERVO5\_FUNCTION} & 10 (\emph{CameraTrigger}, MAP2 al pad M5) \\
        \texttt{CAM1\_TYPE} & 1 (Servo) \\
        \texttt{CAM1\_DURATION} & 1.0 (en \textbf{segons}, no ms): el MAP2 en mode IS-T necessita ordres d'1~s \\
        \texttt{CAM1\_SERVO\_ON} & 1700 (mode IS-T) \\
        \texttt{CAM1\_SERVO\_OFF} & 1500 (dins el rang de repòs del MAP2, 1400--1600) \\
        \texttt{SERVO5\_TRIM} & 1500 \\
        \texttt{CAM1\_TRIGG\_DIST} & 0 a la configuració base; es fixa a cada missió (\cref{sec:planificacio-vol}) \\
        \texttt{SERVO6\_FUNCTION} & $-1$ (GPIO per al \emph{feedback} del SYNC2, pad M6) \\
        \texttt{CAM1\_FEEDBAK\_PIN} & 55 \\
        \texttt{CAM1\_FEEDBAK\_POL} & 1 \\
        \bottomrule
    \end{tabular}
\end{table}

El nom del paràmetre és \texttt{CAM1\_FEEDBAK\_PIN}, escrit \textbf{sense la ``C''} de \emph{feedback} (nomenclatura històrica d'ArduPilot). El pin 55 correspon al pad de \texttt{SERVO6} reassignat com a GPIO.

Amb un objectiu manual sense contactes electrònics, el MAP2 no ha de treballar en mode AF-T (pre-enfocament d'1~s inútil, que afegeix retard constant). El detall del diagnòstic és a \cref{subsec:problemes-camera}.

\subsection{Ajust de la càmera}
\label{subsec:ajust-camera}

\begin{table}[H]
    \centering
    \caption{Ajust de la Sony $\alpha$6000 per a fotogrametria}
    \label{tab:ajust-camera}
    \begin{tabular}{ll}
        \toprule
        Ajust & Valor \\
        \midrule
        Mode d'exposició & M (manual) \\
        Temps d'obturació & 1/1600~s \\
        Diafragma & f/5,6 (anell fixat amb cinta) \\
        Enfocament & Manual a infinit (els dos anells fixats amb cinta) \\
        ISO & Automàtic, amb sostre a 1600 \\
        Mode de dispar & \emph{Single shot}, JPEG Fine \\
        Obturador & Electrònic de primera cortina activat \\
        \emph{Release w/o Lens} & Activat (necessari amb objectiu sense contactes) \\
        \bottomrule
    \end{tabular}
\end{table}

L'òptica manual no transmet dades a la càmera: l'EXIF grava \texttt{FocalLength=0} i \texttt{LensModel} buit, i l'orientació dependria del rumb (\cref{sec:processament-fotos}). Aquestes metadades es reescriuen en el post-procés.

La càmera no té marge dinàmic per compensar la poca llum: les fotografies s'han de fer amb el sol a més de 10--15$^{\circ}$ sobre l'horitzó (\cref{sec:vol6}).

\subsection{Retard de l'obturador}

L'ordre de dispar del controlador (\texttt{TRIG}) i l'exposició real estan separats per $\approx 0{,}28$~s (\cref{sec:retard-obturador}); a 5~m/s són $\approx 1{,}4$~m. Per això el geoetiquetat fa servir la posició del missatge \texttt{CAM} (instant del \emph{feedback} de la sabata) quan està disponible, i es compensa amb $v\cdot\tau$ quan no ho està.

\section{Configuració GPS \emph{moving baseline} i \emph{yaw}}
\label{subsec:gps-moving-baseline}

\begin{table}[H]
    \centering
    \caption{Paràmetres de \emph{moving baseline} i GPS-for-yaw}
    \label{tab:params-gps-mb}
    \begin{tabular}{ll}
        \toprule
        Paràmetre & Valor \\
        \midrule
        \texttt{GPS1\_TYPE} & 17 (MB-Base) \\
        \texttt{GPS2\_TYPE} & 18 (MB-Rover) \\
        \texttt{GPS\_AUTO\_CONFIG} & 1 \\
        \texttt{GPS\_DRV\_OPTIONS} & 0 \\
        \texttt{GPS\_INJECT\_TO} & 127 \\
        \texttt{EK3\_SRC1\_POSXY} / \texttt{VELXY} & 3 (GPS) \\
        \texttt{EK3\_SRC1\_YAW} & 2 (GPS yaw) \\
        \texttt{COMPASS\_ENABLE} & 0 (sense brúixola) \\
        \texttt{GPS1\_POS\_X/Y/Z} & $0$ / $-0{,}19$ / $-0{,}10$~m \\
        \texttt{GPS2\_POS\_X/Y/Z} & $0$ / $+0{,}19$ / $-0{,}10$~m \\
        \texttt{GPS1\_MB\_OFS\_Y} & $-0{,}38$~m \\
        \bottomrule
    \end{tabular}
\end{table}

\texttt{GPS\_DRV\_OPTIONS} ha de ser 0: el valor 1 (ús de l'UART2 per al \emph{moving baseline}) assumeix un cable directe entre els UART2 dels dos F9P per transportar el flux RTCM, que en aquest muntatge no existeix; amb el bit activat el \emph{rover} no arriba mai a RTK Fixed (\texttt{PreArm: GPS[2] yaw not available}). El flux entre receptors passa pel controlador de vol.

\section{Control, bateria i seguretat}
\label{sec:params-control}

\begin{table}[H]
    \centering
    \caption{Paràmetres de control, bateria i seguretat de la configuració de referència}
    \label{tab:params-control}
    \small
    \begin{tabular}{ll}
        \toprule
        Paràmetre & Valor \\
        \midrule
        \texttt{FRAME\_CLASS} & 1 (Quad) \\
        \texttt{FRAME\_TYPE} & 1 (X) \\
        \texttt{MOT\_PWM\_TYPE} & 5 (DShot300)~\cite{ardupilot_dshot} \\
        \texttt{ATC\_RAT\_RLL\_P} / \texttt{ATC\_RAT\_PIT\_P} & 0,07 \\
        \texttt{ATC\_RAT\_RLL\_I} / \texttt{ATC\_RAT\_PIT\_I} & 0,07 \\
        \texttt{ATC\_RAT\_RLL\_D} / \texttt{ATC\_RAT\_PIT\_D} & 0,002 \\
        \texttt{ATC\_RAT\_YAW\_P} & 0,10 \\
        \texttt{ATC\_RAT\_YAW\_I} & 0,01 \\
        \texttt{MOT\_THST\_EXPO} & 0,65 \\
        \texttt{MOT\_SPIN\_ARM} & 0,04 \\
        \texttt{MOT\_SPIN\_MIN} & 0,06 \\
        \texttt{MOT\_THST\_HOVER} & 0,225 (après en vol, \texttt{MOT\_HOVER\_LEARN=2}) \\
        \texttt{BATT\_MONITOR} & 4 (tensió i corrent analògics) \\
        \texttt{BATT\_VOLT\_MULT} & 10,82 (calibrat) \\
        \texttt{BATT\_AMP\_PERVLT} & 87,73 (calibrat) \\
        \texttt{BATT\_LOW\_VOLT} / \texttt{BATT\_FS\_LOW\_ACT} & 21,6~V / 2 (RTL) \\
        \texttt{BATT\_CRT\_VOLT} / \texttt{BATT\_FS\_CRT\_ACT} & 20,4~V / 1 (Land) \\
        \texttt{BATT\_FS\_VOLTSRC} & 0 (tensió instantània) \\
        \texttt{FENCE\_ENABLE} / \texttt{FENCE\_TYPE} & 1 / 7 (altura, cercle i polígon) \\
        \texttt{FENCE\_RADIUS} & 700~m \\
        \texttt{FENCE\_ALT\_MAX} & 120~m \\
        \texttt{WPNAV\_SPEED} & 500~cm/s (5~m/s) en missions automàtiques \\
        \texttt{FLTMODE1} a \texttt{FLTMODE6} & Stabilize, AltHold, Loiter, Auto, Guided, RTL \\
        \texttt{LOG\_BITMASK} & 196607 \\
        \bottomrule
    \end{tabular}
\end{table}

\section{Checklist pre-vol}
\label{sec:checklist-prevol}

\begin{enumerate}[label=\arabic*.]
    \item Verificar que els dos F9P (base i dron) reporten \texttt{RTK Fixed} abans d'armar, amb tipus de \emph{fix} correcte a tots dos receptors (no \texttt{RTK Float} ni \texttt{3D Fix}).
    \item Verificar que el \emph{yaw} GNSS (\emph{moving baseline}) és vàlid i coherent amb el \emph{heading} real (marge $< 1^{\circ}$).
    \item Comprovar l'enllaç RFD868x GS$\leftrightarrow$dron i la recepció de correccions RTCM al dron.
    \item Comprovar el nivell de bateria del dron (tensió per cel·la) i de la base (UPS).
    \item Comprovar la targeta SD i la bateria de la càmera, i que l'ajust és el de \cref{tab:ajust-camera}.
    \item Provar el disparador amb \emph{Trigger Camera Now} i confirmar que apareixen els missatges \texttt{CAM} (\emph{feedback} SYNC2 actiu).
    \item Revisar els modes de vol assignats i que el \emph{geofence} és actiu (\cref{tab:params-control}).
    \item Comprovar a Mission Planner que l'EKF3 és sa (sense advertiments de variància de posició o velocitat), que el HDOP és baix (recomanat $< 1$ amb RTK Fixed) i que no hi ha cap missatge de \emph{prearm} pendent.
    \item Verificar que \texttt{COMPASS\_ENABLE=0} i que no hi ha cap \emph{prearm error} de brúixola.
    \item Comprovar els \emph{failsafes}: pèrdua de RC, pèrdua de telemetria, bateria baixa i crítica.
    \item Comprovar que el sol és a més de 10--15$^{\circ}$ sobre l'horitzó i que l'àrea de vol i l'espai aeri ho permeten (\cref{sec:normativa}).
\end{enumerate}

\section{Planificació de vol}
\label{sec:planificacio-vol}

\subsection{Càlcul del GSD}

El \emph{Ground Sample Distance} depèn només de l'alçada de vol $H$, la mida del sensor, la distància focal $f$ i la resolució de la imatge:
\[
\mathrm{GSD}\,\text{(cm/px)} = \frac{H\,\text{(m)} \times \text{amplada del sensor (mm)} \times 100}{f\,\text{(mm)} \times \text{amplada de la imatge (px)}}
\]

Amb la Sony $\alpha$6000 (sensor APS-C de 23,5~$\times$~15,6~mm, 6000~$\times$~4000~px, mida de píxel $\approx 3{,}92$~\textmu m) i l'òptica TTArtisan de 25~mm:
\[
\mathrm{GSD}\,\text{(cm/px)} \approx 0{,}0157 \times H\,\text{(m)}
\]

\begin{table}[H]
    \centering
    \caption{GSD i cobertura per imatge segons l'alçada (Sony $\alpha$6000 + TTArtisan 25~mm)}
    \label{tab:gsd-altura}
    \begin{tabular}{ccc}
        \toprule
        Alçada (m) & GSD (cm/px) & Empremta (m, ample $\times$ alt) \\
        \midrule
        30 & 0,47 & 28,2 $\times$ 18,7 \\
        50 & 0,78 & 47,0 $\times$ 31,2 \\
        60 & 0,94 & 56,4 $\times$ 37,4 \\
        80 & 1,25 & 75,2 $\times$ 49,9 \\
        100 & 1,57 & 94,0 $\times$ 62,4 \\
        120 & 1,88 & 112,8 $\times$ 74,9 \\
        \bottomrule
    \end{tabular}
\end{table}

Els valors mesurats als vols de prova concorden amb aquesta relació: tenint en compte la focal calibrada pel processat (25,72~mm), un GSD de 0,451~cm/px correspon a 29,6~m d'alçada (vol previst a 30~m), i un GSD de 0,743~cm/px a 48,9~m (vol previst a $\sim$50~m).

\subsection{Paràmetres de vol}

\begin{table}[H]
    \centering
    \caption{Paràmetres de vol adoptats i recomanats}
    \label{tab:parametres-vol}
    \begin{tabular}{lP{9cm}}
        \toprule
        Paràmetre & Valor \\
        \midrule
        Alçada de vol & Segons el GSD objectiu (\cref{tab:gsd-altura}); 30~m i 50~m als vols de prova \\
        Solapament frontal / lateral & 75--77~\% / 65--69~\% al vol 5; 80~\% / 70~\% recomanat amb relleu complex \\
        Velocitat & 5~m/s (\texttt{WPNAV\_SPEED=500}); limitada pel \emph{motion blur}: $v_{\max} = \mathrm{GSD}/t_{\text{exposició}}$ \\
        Dispar & Per distància (\texttt{CAM1\_TRIGG\_DIST}), amb l'interval entre fotos superior a l'1~s de \texttt{CAM1\_DURATION} \\
        Orientació de la càmera & Costat llarg del sensor al llarg del sentit de vol \\
        Il·luminació & Sol a més de 10--15$^{\circ}$ sobre l'horitzó \\
        \bottomrule
    \end{tabular}
\end{table}

Els percentatges de solapament de 80/70~\% són el valor per defecte habitual en programari de fotogrametria per a projectes de precisió amb bona textura del terreny; es poden reduir amb precaució gràcies al geoetiquetat RTK, però un solapament insuficient continua limitant la qualitat del \emph{matching} SfM. Com que \texttt{CAM1\_DURATION=1.0} fixa un interval mínim d'aproximadament 1~s entre dispars, a 5~m/s la distància entre fotos ha de ser d'almenys 5~m (\cref{subsec:autonomia-altura}).

\section{Protocol de camp}
\label{sec:protocol-camp}

\subsection{Posicionar la base}
\begin{enumerate}[label=\alph*)]
    \item Trípode sobre un punt estable amb el cel net a 360$^{\circ}$, lluny de parets, cobertes metàl·liques i vehicles; mesurar l'alçada de l'antena (122~cm als vols de prova).
    \item Engegar la Pi i esperar uns 30~s que aparegui la xarxa \texttt{BaseRTK}; connectar-hi la GS.
    \item Reiniciar el \emph{survey-in} per esborrar la mitjana de la sessió anterior i esperar que sigui vàlid, des del PC connectat a \texttt{BaseRTK} i dins la carpeta \texttt{base\_rtk}: \texttt{python base\_rtk.py survey}. L'ordre va mostrant la durada, les observacions i l'error mitjà, i acaba quan el \emph{survey-in} és vàlid (mínim 300~s i error mitjà, \texttt{meanAcc}, de 2~m o millor) o als 20~min (\texttt{--minuts N} canvia el límit). Mentre dura, el servei \texttt{rtkbase} està aturat i el port 5000 no serveix; es reengega sol al final, també si es talla la connexió SSH o es prem Ctrl-C.
    \item Si no arriba a ser vàlid, la mitjana es queda estancada per \emph{multipath} o per obstrucció del cel: cal millorar la ubicació de l'antena, no allargar l'espera. El receptor continua amb el \emph{survey-in} en marxa, i \texttt{python base\_rtk.py survey --estat} el consulta sense reiniciar-lo.
    \item Quan el \emph{survey-in} és vàlid, l'ordre desa automàticament la posició ECEF (\texttt{meanX/Y/Z}) i la \texttt{meanAcc} finals a \texttt{survey.txt}, dins la sessió activa de la Pi. És l'únic registre de la posició de la base d'aquella sessió, i es baixa amb el registre (vegeu «Registre de la sessió» més avall).
\end{enumerate}

Mentre el \emph{survey-in} no és vàlid, el flux no porta el missatge 1005 (posició de la base) i el dron no pot arribar a RTK Fixed. La posició de la base és, doncs, la del \emph{survey-in}, amb un error de l'ordre de metres que desplaça tota la sessió (\cref{subsec:efecte-surveyin}). Per afirmar posicions absolutes caldrà reutilitzar una posició fixa registrada al mateix punt, o bé corregir el desplaçament de cada sessió amb una solució PPP del seu registre d'observacions.

\paragraph{Registre de la sessió.} El servei \texttt{rtkbase} desa el flux del receptor a fitxers mentre el reenvia pel port 5000, durant tota la sessió: \texttt{str2str} té dues sortides, \texttt{tcpsvr://:5000} i \texttt{file:///home/arnau/logs/actual/base\_\%Y\%m\%d\_\%h\%M.rtcm3::S=1}. Cada arrencada del servei obre una \emph{sessió} nova, una carpeta \texttt{/home/arnau/logs/sNNNN/} (\texttt{actual} apunta a la que està en curs), i dins s'hi crea un fitxer cada hora (\texttt{::S=1}). El nom dels fitxers porta l'hora UTC de la Pi, i la data és poc fiable, perquè la Pi no té rellotge de temps real ni internet: l'hora vàlida és la GNSS de dins de les dades. Un reinici del servei o un tall d'alimentació obre una sessió nova, tret del que fa \texttt{survey}, que continua a la mateixa. El registre és el mateix flux del port, uns 7~MB per hora, amb les observacions RTCM3 (pseudodistància, fase i C/N0 dels quatre sistemes), el missatge 1005 un cop el \emph{survey-in} és vàlid i els missatges UBX de navegació que el receptor emet pel mateix port. No porta efemèrides, que cal agafar d'una estació de referència propera, ni Doppler.

Al final de la sessió es baixa el registre abans d'apagar la base, des del PC connectat a \texttt{BaseRTK}, amb \texttt{python base\_rtk.py baixa} (la darrera sessió) o \texttt{python base\_rtk.py baixa N} (la sessió \texttt{N}; \texttt{python base\_rtk.py llista} mostra les que hi ha, amb mida i hores). L'ordre uneix els fitxers horaris de la sessió en un de sol, \texttt{base\_sNNNN.rtcm3}, en comprova la mida i, si hi ha notes del \emph{survey-in}, baixa també \texttt{base\_sNNNN\_survey.txt}. \texttt{-o carpeta} indica on es deixen, i no s'esborra res de la Pi. Es pot convertir el fitxer a RINEX amb \texttt{convbin -r rtcm3} per calcular a posteriori la posició de la sessió amb una solució PPP o relativa a una estació de referència. Aquest càlcul encara no s'ha fet amb cap registre.

Al disc de la Pi hi caben unes 3\,500~h de registre i no hi ha esborrat automàtic: les sessions que ja no calen s'esborren a mà amb \texttt{python base\_rtk.py esborra N} o \texttt{esborra --abans N}. L'ordre demana confirmació i no esborra mai la sessió activa. Les eines i la seva instal·lació són a la carpeta \texttt{base\_rtk}, amb el seu \texttt{README.md}.

\subsection{Verificar el \emph{moving baseline}}
\begin{enumerate}[label=\alph*)]
    \item Mesurar el vector real entre les dues antenes.
    \item Comprovar el \emph{yaw} GNSS contra el \emph{heading} real (marge $< 1^{\circ}$).
\end{enumerate}

\subsection{Vol}
\begin{enumerate}[label=\alph*)]
    \item Base en \texttt{Fixed} $\rightarrow$ la GS injecta RTCM al dron.
    \item Armar \textbf{només} amb RTK Fixed als dos F9P i \emph{yaw} vàlid.
    \item Volar el \emph{grid} planificat en mode Auto.
    \item Descarregar els \emph{logs} del FC (\texttt{.bin}) i les fotografies de la SD.
\end{enumerate}

\subsection{Post-procés}
\begin{enumerate}[label=\alph*)]
    \item Aparellar cada fotografia amb el seu missatge \texttt{CAM} (o \texttt{TRIG}) del \emph{log} i reescriure l'EXIF (\cref{sec:processament-fotos}).
    \item Generar el fitxer de geoetiquetes amb la precisió per foto.
    \item Processar amb ODX/WebODM (\cref{sec:processament}) i generar ortomosaic, DSM i núvol de punts.
\end{enumerate}

\section{Punts de control terrestre (GCP)}
\label{sec:gcp}

La precisió del producte final només es pot afirmar amb punts mesurats de manera independent. Com que la base fa un \emph{survey-in} a cada sessió, el \textbf{marc de coordenades} dels punts és tan important com la seva precisió. El protocol previst és:
\begin{itemize}
    \item Dianes ben distribuïdes per tota l'àrea, incloent-hi les cantonades i el centre, marcades amb un patró d'alt contrast reconeixible a la fotografia: com a mínim cinc de control i, a part, deu de comprovació (\cref{subsec:pla-precisio}).
    \item Posició de cada diana mesurada amb RTK contra \textbf{la mateixa base i la mateixa sessió} que el vol, amb una precisió prevista d'1--2~cm. Així els punts i els \emph{geotags} comparteixen el desplaçament del \emph{survey-in} i aquest s'anul·la a la comparació. Si es mesuren amb estació total o amb una altra base, cal ajustar-hi una transformació abans de comparar; si no, l'RMSE el domina l'error de metres de la base.
    \item Separar punts de control (que s'usen a l'ajust) de punts de comprovació (que no hi intervenen), i calcular l'RMSE H/V d'aquests últims.
\end{itemize}

\pendent{Encara no s'ha desplegat cap diana ni mesurat cap punt de control: aquest és el protocol previst, no un resultat. És la tasca que desbloqueja la precisió absoluta.}

\section{Processament (\emph{Structure from Motion} i PPK)}
\label{sec:processament}

El processament fotogramètric es fa amb el motor ODX 3.8.3 (WebODM 3.3.0, NodeODX 2.3.1) en un PC amb un i5-10400F, 32~GB de RAM i una RTX~2060. Un vol de 118 fotografies triga uns 42~minuts. Les opcions que s'aparten dels valors per defecte són: resolució d'ortofoto de 1~cm/px, resolució de DEM de 2~cm/px, qualitat del núvol de punts alta, generació de DSM, \emph{auto-boundary}, COG, glTF i \emph{point cloud} EPT. El detall de cada processat és a \cref{chap:vols-prova}.

Com a alternativa o complement al RTK en temps real, es contempla el post-procés \textbf{PPK} amb RTKLIB a partir dels enregistraments \texttt{UBX} en brut de la base i del dron, que dona una solució de posicionament independent de la qualitat de l'enllaç RTK durant el vol.

\pendent{El PPK no s'ha aplicat encara a cap vol: els processats actuals fan servir les posicions RTK en temps real del controlador.}

\section{Control de qualitat i criteris de validació}
\label{sec:validacio}

\subsection{Efecte del \emph{survey-in} sobre la validació}
\label{subsec:efecte-surveyin}

Amb el \emph{survey-in} la base calcula la seva posició com a mitjana de posicions autònomes, sense correccions, i la difon al missatge 1005. L'error d'aquesta posició es transmet íntegrament al \emph{rover}~\cite{ublox_f9p_im}: tota la sessió queda desplaçada pel mateix vector, que és de l'ordre de metres (el criteri d'acceptació del protocol és una \texttt{meanAcc} d'$\leq 2$~m). En canvi, la posició \emph{relativa} respecte a la base gairebé no es veu afectada: l'error de longitud de la línia de base és el de la base multiplicat per la línia de base i dividit per la distància als satèl·lits ($\sim 2\cdot 10^{7}$~m), uns 0,02~mm amb una base desplaçada 2~m i una línia de base de 200~m.

\begin{itemize}
    \item \textbf{No afectat} (es pot validar sense referència absoluta): l'exactitud de la línia de base RTK, la repetibilitat del \emph{rover}, el \emph{heading} i la geometria interna del producte (forma i escala).
    \item \textbf{Afectat:} la posició absoluta del mapa, la comparabilitat entre sessions i qualsevol comparació amb una referència externa. Com que cada sessió fa un \emph{survey-in} nou, dos mapes del mateix lloc queden desplaçats entre ells de l'ordre de metres. L'alçada de l'antena de la base (122~cm) entra a la cota amb pes 1:1.
\end{itemize}

Per recuperar la posició absoluta hi ha tres vies, no excloents: reutilitzar una posició fixa registrada al mateix punt; mesurar un punt conegut amb la mateixa base i corregir-ne el desplaçament; o calcular la posició de la sessió per PPP a partir del registre de la base i aplicar-ne el vector de correcció a tot el producte. Aquesta última és viable en principi amb el flux RTCM ja generat, que conté observacions MSM (1074, 1084, 1094, 1124) convertibles a RINEX amb RTKLIB (\texttt{convbin}); cal verificar-ho i que la sessió duri prou hores per a un PPP de nivell centimètric.

\subsection{Pla de proves de precisió}
\label{subsec:pla-precisio}

\begin{table}[H]
    \centering
    \caption{Pla de proves de precisió amb \emph{survey-in} i estat actual}
    \label{tab:validacio}
    \footnotesize
    \begin{tabular}{P{0.8cm}P{4.9cm}P{2.8cm}P{2.6cm}P{2.4cm}}
        \toprule
        Prova & Procediment & Marc de referència & Criteri & Estat \\
        \midrule
        T1 & Registrar \texttt{NAV-SVIN} (\texttt{meanX/Y/Z}, \texttt{meanAcc}) a cada sessió i, a mostres, comparar-lo amb la solució PPP del mateix lloc & --- & Informatiu: documenta el desplaçament & \textbf{Pendent}; no registrat als vols 5 i 6 \\
        T2 & \emph{Rover} estàtic a tres distàncies (p.~ex.\ 10, 50 i 100~m): vector base$\rightarrow$\emph{rover} RTK contra una mesura independent (estació total o distanciòmetre i nivell) & Relatiu a la base (independent del \emph{survey-in}) & $< 2$~cm H, $< 3$~cm V & \textbf{Pendent} \\
        T3 & Repetibilitat estàtica: 10--15~min sobre un punt; desviació estàndard H/V i percentatge de RTK Fixed & Relatiu & Coherent amb 0,01~m + 1~ppm & \textbf{Pendent} \\
        T4 & \emph{Heading}: \emph{yaw} GNSS contra un rumb de referència mesurat amb estació total, en diverses orientacions & Independent de la base & $\sigma < 0{,}5^{\circ}$ & \textbf{Pendent}; \emph{yaw} vàlid en tots els vols \\
        T5 & Producte fotogramètric: RMSE dels punts de comprovació, amb i sense punts de control, informant del nombre de punts & Mateix marc que els \emph{geotags} (mateixa base i sessió) & $< 3$~cm H, $< 5$~cm V & \textbf{Pendent}; consistència interna: reprojecció 0,90~px, model--GPS CE90 0,59~m \\
        T6 & Precisió absoluta: aplicar el vector PPP de la base i comprovar sobre una marca geodèsica coneguda & Marc extern (EPSG:32631) & Per definir (limitat pel PPP) & \textbf{Pendent}; requereix el registre de la base \\
        \bottomrule
    \end{tabular}
\end{table}

La comparativa amb un \emph{dataset} de referència DJI P4 RTK només té sentit després de T6, perquè necessita un mapa georeferenciat en absolut.

Sobre el nombre de punts: amb només cinc la incertesa de l'RMSE és gran. L'edició 2 de l'estàndard ASPRS demana un mínim de 30 punts de comprovació per afirmar una precisió de producte, i admet un nombre menor en projectes petits amb UAS si l'informe ho declara~\cite{asprs2023}. Per això es recomana un mínim de deu de comprovació, a més dels de control, i indicar-ne sempre el nombre.

Els indicadors de consistència interna (error de reprojecció, diferència entre el model i les posicions GPS de les fotos) mesuren la coherència del processat, \textbf{no la precisió}: no substitueixen T2--T6. Fins que no es facin, no s'afirma cap precisió absoluta ni centimètrica.

\section{Exportació de resultats}
\label{sec:exportacio}

Cada processat lliura, en el sistema de coordenades EPSG:32631 (UTM 31N), els productes següents: ortofoto GeoTIFF (1~cm/px a 30~m), DSM GeoTIFF (2~cm/px), núvol de punts dens (LAZ i EPT), models texturats (OBJ i GLB) i un informe PDF amb les estadístiques del processat.

\section{Manteniment del dron}
\label{sec:manteniment}

\subsection{Abans de cada vol}
\begin{itemize}
    \item Inspecció visual del \emph{frame}, els braços i els \emph{standoffs}: fissures o joc mecànic.
    \item Hèlixs sense esquerdes ni deformacions, equilibrades i amb els collarins ben apretats.
    \item Connectors XT60/XT90 i JST-GH sense joc ni signes d'escalfament.
    \item Antenes helicals ben fixades, verticals i sense danys al \emph{bracket} de \emph{baseline}.
    \item Estat de la bateria (inflor, tensió per cel·la) i càrrega suficient per al pla de vol.
    \item Targeta SD de la càmera amb espai lliure suficient i formatada correctament.
\end{itemize}

\subsection{Manteniment periòdic}
\begin{itemize}
    \item Revisió dels parells de collament de motors i braços cada poques sortides.
    \item Comprovació de l'estat dels coixinets dels motors (soroll, joc a l'eix).
    \item Actualització del \emph{firmware} d'ArduPilot i dels receptors F9P només de manera controlada, verificant la compatibilitat de paràmetres abans i després.
    \item Verificació periòdica del \emph{heading} GNSS contra una referència coneguda, per detectar deriva mecànica del \emph{bracket}.
    \item Neteja de contactes i connectors, especialment després de volar en ambients humits o amb pols.
    \item Cicles de manteniment de les bateries LiPo: emmagatzematge a tensió de repòs, sense descarregar per sota del llindar del fabricant, i retirada de qualsevol \emph{pack} que hagi patit un \emph{crash}.
\end{itemize}

\section{Normativa i seguretat}
\label{sec:normativa}

Aquesta secció és informativa i no constitueix assessorament legal. El marc normatiu EASA per a drones autoconstruïts està en transició; cal verificar-ne l'estat vigent abans de cada campanya de vol~\cite{easa_open_category}.

\subsection{Classificació del dron}

Aquest UAV és \textbf{autoconstruït} (\emph{privately built}): no prové d'un \emph{kit} llest per muntar comercialitzat com a producte únic, sinó de components individuals. Segons EASA, un dron autoconstruït \textbf{no pot portar marcat de classe C0--C4}, i per tant queda limitat, dins la categoria \emph{Open}, a la subcategoria \textbf{A3} (lluny de persones i de zones edificades), amb pes màxim de 25~kg~\cite{easa_open_category}.

Requisits pràctics derivats:
\begin{itemize}
    \item Registre com a operador de drones a l'autoritat aeronàutica competent (a l'Estat espanyol, AESA) i identificació visible del número de registre a l'aeronau.
    \item Certificat de pilot A1/A3 (formació en línia i examen) com a mínim; l'operador el té.
    \item Vol exclusivament en subcategoria A3: distància horitzontal mínima de 150~m respecte a zones residencials, comercials, industrials o recreatives, i lluny de persones alienes a l'operació.
    \item Comprovació de l'espai aeri i de les restriccions locals (zones \emph{geo-awareness}, NOTAM, proximitat a aeroports) abans de cada vol.
    \item Assegurança de responsabilitat civil, recomanable (obligatòria en molts casos segons el pes i l'ús).
\end{itemize}

Sota les regles transitòries, els drones autoconstruïts sense marcat de classe poden operar segons el règim heretat (per pes) fins a l'1 de gener de 2027. EASA desenvolupa un règim permanent per a drones autoconstruïts, equivalent aproximadament a la classe C4 (subcategoria A3, $<25$~kg, pilot amb certificat A1/A3)~\cite{easa_open_category}.

\subsection{Seguretat operacional}

\begin{itemize}
    \item No volar amb RTK Float ni sense \emph{fix} GNSS vàlid si la missió depèn de precisió centimètrica (a més, redueix la qualitat de l'estimació de posició de l'EKF3).
    \item Definir sempre un \emph{geofence} i accions de \emph{failsafe} per a bateria baixa i pèrdua d'enllaç.
    \item Mantenir contacte visual (VLOS) llevat que es disposi de l'autorització corresponent per a BVLOS.
    \item Avaluació de riscos bàsica de la zona de vol: obstacles, línies elèctriques, tercers i meteorologia (vent, pluja).
\end{itemize}

\section{Línies de continuïtat}
\label{sec:treball-futur}

\begin{itemize}
    \item Validació absoluta amb punts de control i de comprovació (\cref{sec:gcp}).
    \item Millora de l'autonomia: canvi del grup propulsor i bateria Li-ion (\cref{sec:autonomia,sec:propulsor}).
    \item Configuració del filtre \emph{notch} harmònic i Autotune.
    \item Migració a UM982: un sol mòdul amb \emph{heading} natiu, menys cablejat i menys consum.
    \item Integració de LiDAR (Livox Mid-360) per a vegetació densa.
    \item Càmera multiespectral per a agricultura de precisió.
    \item PCB \emph{carrier} pròpia (F9P + SBC + UPS integrats).
    \item Publicar correccions a una xarxa NTRIP comunitària (Centipède~\cite{centipede} o rtk2go).
\end{itemize}

\section{Projectes de referència}
\label{sec:referencies}

\begin{itemize}
    \item \texttt{rtkbase} (Stefal) -- programari de base~\cite{rtkbase}.
    \item SparkFun RTK Express -- disseny KiCad i \emph{firmware}~\cite{sparkfun_rtk_express}.
    \item \texttt{jacobjhansen/RTCM-Pi} -- flux de treball PPP~\cite{rtcm_pi}.
    \item Centipède -- xarxa NTRIP comunitària a Europa~\cite{centipede}.
    \item ArduPilot -- \emph{moving baseline} i GPS-for-yaw~\cite{ardupilot_moving_baseline}.
\end{itemize}
